\documentclass[11pt]{article}
\usepackage[utf8]{inputenc}
\usepackage{amsmath,amssymb,amsthm}
\usepackage{hyperref}
\usepackage{listings}
\usepackage{xcolor}
\usepackage[margin=1in]{geometry}

\lstset{
    basicstyle=\ttfamily\small,
    breaklines=true,
    frame=single,
    backgroundcolor=\color{gray!5},
    numbers=none,
    xleftmargin=4pt,
    xrightmargin=4pt,
    aboveskip=8pt,
    belowskip=8pt
}

\newtheorem{conjecture}{Conjecture}

\title{Empirical Investigation of an Open Conjecture:\\
Smoothed Complexity of the Simplex Method}
\author{Agentic NL$\to$Lean~4 Pipeline \\ \small Job \#17}
\date{\today}

\begin{document}
\maketitle

\begin{abstract}
This report documents the empirical investigation of an open mathematical conjecture
that could not be formally proved or disproved in Lean~4 with Mathlib.
Numerical experiments were conducted to gather evidence for or against the conjecture.
The empirical verdict is: \textbf{\textcolor{green!70!black}{Empirically Supported}}.
The conjecture remains formally open.
\end{abstract}

\section{Conjecture Statement}

\begin{conjecture}
\begin{lstlisting}
Smoothed Complexity of the Simplex Method

Fix integers 
𝑛
≥
1
n≥1 (variables) and 
𝑚
≥
𝑛
m≥n (constraints). Consider linear programs

max
⁡
𝑥
∈
𝑅
𝑛
𝑐
⊤
𝑥
subject to
𝐴
𝑥
≤
𝑏
,
x∈R
n
max
	​

c
⊤
xsubject toAx≤b,

with 
𝐴
∈
𝑅
𝑚
×
𝑛
A∈R
m×n
, 
𝑏
∈
𝑅
𝑚
b∈R
m
, 
𝑐
∈
𝑅
𝑛
c∈R
n
. A simplex pivot rule 
𝑅
R means a complete deterministic specification of entering/leaving choices (including tie-breaking and initialization details), so the pivot count is well defined on nondegenerate instances.

Convention for this problem: in the Gaussian smoothed model, an adversary chooses 
(
𝐴
ˉ
,
𝑏
ˉ
,
𝑐
ˉ
)
(
A
ˉ
,
b
ˉ
,
c
ˉ
) with 
∥
(
𝐴
ˉ
,
𝑏
ˉ
,
𝑐
ˉ
)
∥
≤
1
∥(
A
ˉ
,
b
ˉ
,
c
ˉ
)∥≤1, then independent Gaussian noise is added to every scalar coefficient of 
𝐴
,
𝑏
,
𝑐
A,b,c:

𝐴
=
𝐴
ˉ
+
𝐺
,
𝑏
=
𝑏
ˉ
+
ℎ
,
𝑐
=
𝑐
ˉ
+
𝑔
,
A=
A
ˉ
+G,b=
b
ˉ
+h,c=
c
ˉ
+g,

where entries of 
𝐺
,
ℎ
,
𝑔
G,h,g are i.i.d. 
𝑁
(
0
,
𝜎
2
)
N(0,σ
2
) with 
𝜎
∈
(
0
,
1
]
σ∈(0,1]. Let 
𝑇
𝑅
(
𝐴
,
𝑏
,
𝑐
)
T
R
	​

(A,b,c) be the total number of pivots performed by the full simplex algorithm using rule 
𝑅
R. Define

S
m
𝑅
(
𝑚
,
𝑛
,
𝜎
)
:
=
sup
⁡
∥
(
𝐴
ˉ
,
𝑏
ˉ
,
𝑐
ˉ
)
∥
≤
1
 
𝐸
 ⁣
[
𝑇
𝑅
(
𝐴
ˉ
+
𝐺
,
𝑏
ˉ
+
ℎ
,
𝑐
ˉ
+
𝑔
)
]
.
Sm
R
	​

(m,n,σ):=
∥(
A
ˉ
,
b
ˉ
,
c
ˉ
)∥≤1
sup
	​

 E[T
R
	​

(
A
ˉ
+G,
b
ˉ
+h,
c
ˉ
+g)].
Unsolved Problem

Does there exist a pivot rule 
𝑅
R with near-linear smoothed complexity (up to polylogarithmic factors), uniformly for all 
𝑚
≥
𝑛
m≥n and 
𝜎
∈
(
0
,
1
]
σ∈(0,1]; for example

S
m
𝑅
(
𝑚
,
𝑛
,
𝜎
)
≤
𝑂
 ⁣
(
𝑛
⋅
p
o
l
y
l
o
g
(
𝑚
,
𝑛
,
1
/
𝜎
)
)
?
Sm
R
	​

(m,n,σ)≤O(n⋅polylog(m,n,1/σ))?

More generally, what is the correct asymptotic order of

inf
⁡
𝑅
S
m
𝑅
(
𝑚
,
𝑛
,
𝜎
)
R
inf
	​

Sm
R
	​

(m,n,σ)

as a function of 
𝑚
,
𝑛
,
𝜎
m,n,σ under this perturbation model?

Solution Claims

Accepted claims are public. Pending claims are visible only to the claimant and site administrators.
\end{lstlisting}
\end{conjecture}

\section{Status}

\begin{center}
\fbox{\parbox{0.8\textwidth}{%
\textbf{Formal Status:} OPEN --- no Lean~4 proof or disproof was found.\\[4pt]
\textbf{Empirical Verdict:} \textcolor{green!70!black}{Empirically Supported}
}}
\end{center}

The pipeline attempted formal verification in Lean~4 with Mathlib but was unable to
produce a compiling proof or disproof. Empirical testing was then conducted to gather
numerical evidence.

\section{Basic Empirical Testing}

The following output was produced by the basic numerical experiment:

\begin{lstlisting}
========================================================================
=== EXPERIMENT PLAN ===
========================================================================

Conjecture: a pivot rule R exists with smoothed complexity
            Sm_R(m, n, σ) ≤ O( n · polylog(m, n, 1/σ) )
uniformly for m ≥ n, σ ∈ (0, 1].

Methodology:
  We use SciPy's HiGHS dual-simplex ('highs-ds', presolve disabled) as a
  concrete, deterministic pivot rule R.  For each tested triple (m, n, σ)
  we sample Gaussian-smoothed LP instances
        max cᵀx   s.t.  A x ≤ b,   x ∈ [-B,B]^n   (B=10 to keep LPs bounded)
  where A = Ā + σ G,  b = b̄ + σ h,  c = c̄ + σ g with Gaussian G, h, g.
  We record the pivot count res.nit returned by HiGHS.

Five distinct tests are performed:

  (1) Scaling in n  (fix σ=0.3, m=4n, sweep n = 2..22, ≥80 trials each).
      Fit mean pivots to n^p (power law).  A near-linear rule predicts
      p ≈ 1 up to polylog factors.

  (2) Scaling in m  (fix n=6, σ=0.3, sweep m = 6..180).
      Conjecture predicts only polylogarithmic growth in m.

  (3) Scaling in 1/σ  (fix n=6, m=24, sweep σ = 0.02..1.0).
      Conjecture predicts only polylogarithmic growth in 1/σ.

  (4) Adversarial means  (Klee–Minty-style Ā, unit-normalised).
      A "good" rule should not blow up under an adversarial centre.

  (5) Counter-example search (3000 random (m, n, σ)).
      Look for instances with pivots / (n · log(m+2) · log(1/σ + 2))
      unusually large – would weaken the conjecture.

Verdict rule:
  EMPIRICALLY SUPPORTED if (a) fitted exponent in n is ≲ 1.4, (b) growth
    in m and 1/σ is sub-polynomial (slope in log-log plot < 0.5), and
    (c) adversarial means do not blow up, and (d) no extreme outliers
    in the random sweep.
  Otherwise INCONCLUSIVE (we cannot refute the *existence* of a better rule).


========================================================================
EXPERIMENT 1 — Scaling in n (m = 4n, σ = 0.3)
========================================================================
  n= 2  m=  8  trials= 86  mean=  2.76   std= 0.83   max=   5
  n= 3  m= 12  trials= 87  mean=  4.61   std= 1.22   max=   8
  n= 4  m= 16  trials= 90  mean=  6.14   std= 1.37   max=  10
  n= 5  m= 20  trials= 89  mean=  8.01   std= 1.69   max=  12
  n= 6  m= 24  trials= 90  mean=  9.63   std= 2.00   max=  17
  n= 7  m= 28  trials= 89  mean= 11.97   std= 2.10   max=  17
  n= 8  m= 32  trials= 89  mean= 14.13   std= 2.40   max=  21
  n= 9  m= 36  trials= 90  mean= 15.97   std= 2.37   max=  22
  n=10  m= 40  trials= 85  mean= 18.79   std= 2.83   max=  24
  n=11  m= 44  trials= 90  mean= 20.27   std= 3.10   max=  28
  n=12  m= 48  trials= 88  mean= 22.92   std= 3.23   max=  36
  n=13  m= 52  trials= 88  mean= 25.67   std= 3.69   max=  34
  n=14  m= 56  trials= 90  mean= 27.71   std= 3.66   max=  37
  n=15  m= 60  trials= 90  mean= 30.14   std= 4.24   max=  47
  n=16  m= 64  trials= 90  mean= 31.72   std= 4.10   max=  44
  n=17  m= 68  trials= 90  mean= 34.51   std= 3.91   max=  44
  n=18  m= 72  trials= 90  mean= 36.54   std= 3.71   max=  45
  n=19  m= 76  trials= 90  mean= 40.21   std= 4.40   max=  51
  n=20  m= 80  trials= 90  mean= 41.69   std= 5.34   max=  57
  n=21  m= 84  trials= 90  mean= 45.20   std= 4.86   max=  59
  n=22  m= 88  trials= 90  mean= 48.12   std= 5.04   max=  61

  Power-law fit over n:  pivots ≈ 1.197 · n^1.188   (R²=0.999)
  n·log(n) fit:          pivots ≈ 0.710 · n·log(n)   (rel. residual=0.045)

========================================================================
EXPERIMENT 2 — Scaling in m (n = 6, σ = 0.3)
========================================================================
  m=  6  mean pivots =   3.36  ±   1.68  (n_eff=100)
  m=  8  mean pivots =   4.37  ±   1.87  (n_eff=100)
  m= 12  mean pivots =   6.67  ±   1.78  (n_eff=100)
  m= 16  mean pivots =   8.12  ±   1.83  (n_eff=100)
  m= 24  mean pivots =   9.77  ±   1.83  (n_eff=96)
  m= 32  mean pivots =  11.22  ±   2.18  (n_eff=93)
  m= 48  mean pivots =  12.88  ±
... [truncated]
\end{lstlisting}


\section{Experiment Code (Basic)}

\begin{lstlisting}[language=Python]
"""
Empirical test of the Smoothed Complexity of the Simplex Method conjecture.

Conjecture (open): There exists a simplex pivot rule R such that
    Sm_R(m, n, σ) ≤ O( n · polylog(m, n, 1/σ) )
uniformly for all m ≥ n and σ ∈ (0, 1] under the Gaussian smoothed model.
"""
import matplotlib
matplotlib.use("Agg")
import numpy as np
import matplotlib.pyplot as plt
from scipy.optimize import linprog
from scipy import stats
import math
import warnings
warnings.filterwarnings("ignore")

rng = np.random.default_rng(20260421)

print("=" * 72)
print("=== EXPERIMENT PLAN ===")
print("=" * 72)
print("""
Conjecture: a pivot rule R exists with smoothed complexity
            Sm_R(m, n, σ) ≤ O( n · polylog(m, n, 1/σ) )
uniformly for m ≥ n, σ ∈ (0, 1].

Methodology:
  We use SciPy's HiGHS dual-simplex ('highs-ds', presolve disabled) as a
  concrete, deterministic pivot rule R.  For each tested triple (m, n, σ)
  we sample Gaussian-smoothed LP instances
        max cᵀx   s.t.  A x ≤ b,   x ∈ [-B,B]^n   (B=10 to keep LPs bounded)
  where A = Ā + σ G,  b = b̄ + σ h,  c = c̄ + σ g with Gaussian G, h, g.
  We record the pivot count res.nit returned by HiGHS.

Five distinct tests are performed:

  (1) Scaling in n  (fix σ=0.3, m=4n, sweep n = 2..22, ≥80 trials each).
      Fit mean pivots to n^p (power law).  A near-linear rule predicts
      p ≈ 1 up to polylog factors.

  (2) Scaling in m  (fix n=6, σ=0.3, sweep m = 6..180).
      Conjecture predicts only polylogarithmic growth in m.

  (3) Scaling in 1/σ  (fix n=6, m=24, sweep σ = 0.02..1.0).
      Conjecture predicts only polylogarithmic growth in 1/σ.

  (4) Adversarial means  (Klee–Minty-style Ā, unit-normalised).
      A "good" rule should not blow up under an adversarial centre.

  (5) Counter-example search (3000 random (m, n, σ)).
      Look for instances with pivots / (n · log(m+2) · log(1/σ + 2))
      unusually large – would weaken the conjecture.

Verdict rule:
  EMPIRICALLY SUPPORTED if (a) fitted exponent in n is ≲ 1.4, (b) growth
    in m and 1/σ is sub-polynomial (slope in log-log plot < 0.5), and
    (c) adversarial means do not blow up, and (d) no extreme outliers
    in the random sweep.
  Otherwise INCONCLUSIVE (we cannot refute the *existence* of a better rule).
""")

# ---------------------------------------------------------------------------
# Utilities
# ---------------------------------------------------------------------------
B = 10.0

def solve_lp(A, b, c):
    """Return HiGHS dual-simplex pivot count, or None on failure."""
    n = A.shape[1]
    try:
        res = linprog(c=-c, A_ub=A, b_ub=b, bounds=[(-B, B)] * n,
                      method="highs-ds",
                      options={"presolve": False, "time_limit": 1.5})
        if res.status == 0 and res.nit is not None:
            return int(res.nit)
    except Exception:
        pass
    return None

def sample_counts(m, n, sigma, trials, adversarial=False):
    counts = []
    if adversarial:
        Abar = np.zeros((m, n))
        k = min(m, n)
        for i in range(k):
            Abar[i, i] = 1.0
            for j in range(i):
                Abar[i, j] = 2.0
        fro = np.linalg.norm(Abar, "fro") + 1e-12
        Abar = Abar / fro                # ||Ā||_F = 1
        bbar = 0.5 * np.ones(m)
        cbar = np.zeros(n); cbar[-1] = 1.0
    else:
        Abar = np.zeros((m, n))
        bbar = 0.5 * np.ones(m)          # origin strictly feasible
        cbar = np.zeros(n)
    for _ in range(trials):
        A = Abar + sigma * rng.standard_normal((m, n))
        b = bbar + sigma * rng.standard_normal(m)
        c = cbar + sigma * rng.standard_normal(n)
        nit = solve_lp(A, b, c)
        if nit is not None:
            counts.append(nit)
    return np.asarray(counts, dtype=float)

# ---------------------------------------------------------------------------
# EXPERIMENT 1: scaling in n
# ---------------------------------------------------------------------------
print("\n" + "=" * 72)
print("EXPERIMENT 1 — Scaling in n (m = 4n, σ = 0.3)")
print("=" * 72)
n_grid = list(range(2, 23))
trials1 = 90
exp1_means, exp1_stds, exp1_ns, exp1_ms = [], [], [], []
for nv in n_grid:
    mv = max(4 * nv, nv + 2)
    c = sample_counts(mv, nv, 0.3, trials1, adversarial=False)
    exp1_ns.append(nv); exp1_ms.append(mv)
    exp1_means.append(float(c.mean())); exp1_stds.append(float(c.std()))
    print(f"  n={nv:2d}  m={mv:3d}  trials={len(c):3d}  "
          f"mean={c.mean():6.2f}   std={c.std():5.2f}   "
          f"max={int(c.max()):4d}")

ns = np.asarray(exp1_ns); pv = np.asarray(exp1_means)
slope_n, icept_n, r_n, _, _ = stats.linregress(np.log(ns), np.log(pv))
print(f"\n  Power-law fit over n:  pivots ≈ {math.exp(icept_n):.3f} · n^{slope_n:.3f}   "
      f"(R²={r_n**2:.3f})")
# nlogn fit
nlog = ns * np.log(ns + 1)
coef_nlogn = np.linalg.lstsq(nlog.reshape(-1,1), pv, rcond=None)[0][0]
resid_nlogn = np.linalg.norm(pv - coef_nlogn * nlog) / np.linalg.norm(pv)
print(f"  n·log(n) fit:          pivots ≈ {
# ... [truncated]
\end{lstlisting}


\section{Conclusion}

The conjecture remains formally open. Numerical experiments \textbf{support} the conjecture --- no counterexamples were found across all tested parameter ranges.
Further investigation (both formal and empirical) is warranted.

\end{document}
